\section{Cross-level accountability framework}

The three aggregation levels are connected by a shared accountability framework:
the same fail-closed discipline is applied to auditing scanner leakage, center
leakage, descriptive task accessibility, region retrieval, whole-slide
aggregation, institutional influence, fairness constraints, robustness,
provenance, and claim validity. This section is the conceptual glue of the
manuscript.

\subsection{Audit endpoints at each level}

\begin{itemize}
\item \textbf{Representation level:} scanner and acquisition recoverability
  (linear and nonlinear probes with paired permutation nulls), descriptive
  tissue-category accessibility under the corrected fixed-estimand, region
  retrieval and same-region cosine similarity, acquisition-category leakage, and
  spectral accessibility.
\item \textbf{Whole-slide level:} slide-level ordinal metrics (QWK), matched
  controlled comparisons across MIL architectures, computational efficiency and
  stability, and branch-contribution attribution.
\item \textbf{Institutional level:} institutional influence weights,
  dominant-site and corruption stress, ordinal-shift detection, held-out-center
  accuracy, privacy and secure-aggregation behavior, and communication
  accounting.
\end{itemize}

\subsection{How artifacts propagate across levels}

Scanner and center shortcuts can propagate from representation features into MIL
aggregation and federated updates. The framework makes the propagation
checkable: a frozen feature array can be audited for scanner and center
recoverability before it enters a whole-slide model or a federated client. An
explicit acquisition or institution branch can be bottlenecked, inspected, and
(in the presence of verified swap metadata) swapped, which a blind linear
projection cannot provide.

\subsection{Claim validity as an audit artifact}

Every active numerical claim binds to an immutable result artifact with a
recorded hash; every architectural claim binds to source and tests; every
theoretical or protocol claim binds to its exact specification file. Negative
results are reported in the main text. Pending controlled validation is labeled
pending. This framework is what allows the manuscript to present a negative
neural-increment result, a pending TransnnMIL superiority evaluation, and a
partially specified institutional-weighting protocol without conflating their
evidence maturity.
